\documentclass[../main.tex]{subfiles}
\begin{document}
%==========================================TIÊU ĐỀ TẬP========================================
\begin{table}[h]
    \begin{adjustwidth}{-1cm}{}
        \begin{tabular}{>{\centering\arraybackslash}p{6cm}|>{\raggedright\arraybackslash}p{12.5cm}}
            \multirow{5}{6cm}
            {\\[-40pt]
                \flaregothic\fontsize{20pt}{20pt}\selectfont \centering
                \color{eptype}HISTOIRE\color{black}\\
                \freshbot\fontsize{72pt}{72pt}\selectfont
                \color{epnum}40\\[6pt]
                \cafeteria\fontsize{20pt}{20pt}\selectfont\color{epnum}(D-40)\color{black}
                \\[18pt]
                \color{black}
            }
             &
            % Dòng thứ nhất: tên tập tiếng Nhật
            {\fotskip\fontsize{18pt}{18pt}\selectfont \ruby{氷}{ひ}\ruby{川}{かわ}\ruby{家}{け}の\ruby{三}{さん}\ruby{匹}{びき}}
            \\[6pt]
            % Dòng thứ hai: tên tập tiếng Anh
             & {\cafeteria\fontsize{18pt}{18pt}\selectfont The Three Dogs of the Hikawa Family}                                         \\[4pt]
            % Dòng thứ ba: tên tập tiếng Việt
             & {\vnmsans\fontsize{18pt}{18pt}\selectfont Ba chú chó nhà Hikawa}                                                         \\[8pt]
            % Dòng thứ tư: Nhân vật xuất hiện
             & {\caviar{\fontsize{13pt}{13pt}\selectfont Nhân vật: \emp{kuon2} \emp{ryuuhou2} \emp{kaiho2} \emp{suzuno2} \emp{naomi2}}}
            \\[8pt]
        \end{tabular}
    \end{adjustwidth}
\end{table}
%=========================================TRUYỆN CHÍNH========================================
\color{black}

\txtbx[2]{
    {\color{vol} \uline{Tóm tắt:}} Giữa nắng gió thoang thoảng của vùng ngoại ô Kawachi vào một buổi chiều tháng Tám, gian phòng thí nghiệm vừa hoàn thiện bên bãi đất hoang bỗng rộn ràng bởi những bước chân lúp xúp của Tô và Ki. Hai chú cún là những mảnh ghép sống động gợi nhớ về Mi, cô chó trung thành đã gắn bó mười bốn năm cùng đại gia đình Hikawa. Từ khoảnh khắc cô bé Ryuuhou hai tuổi nằng nặc đòi chọn nàng cún Bắc Kinh lai Nhật tại tiệm thú cưng, qua những năm tháng sóng gió ly thân đầy cô quạnh mà hai bố con Kuon cùng Mi nương tựa chăm sóc lẫn nhau, cho tới hai lần vượt cạn ở tuổi xế chiều và giấc ngủ thiên thu dưới tán rẻ quạt đầu tháng Sáu, hình bóng Mi vẫn vẹn nguyên như một người bảo hộ thầm lặng che chở cho mái ấm Kawachi.
}

Gió chiều cuối tháng Bảy thổi qua khoảng sân rộng phía sau xưởng cơ khí của bố Ryujin, mang theo mùi dầu máy quen thuộc quyện cùng hương khoáng thạch mát lạnh từ tàn tích nhà máy đá cũ và vị ngòn ngọt của thảm cỏ dại đang tắm mình dưới nắng xế chiều.

Ngay bên cạnh bức tường rào ranh giới ngăn cách xưởng cơ khí với khu đất hoang, cánh cửa kính chống cháy của gian phòng thí nghiệm mới tinh của Suzuno đang khẽ hé mở. Một làn gió nhẹ luồn qua ô thông gió hiện đại, xua tan mùi sơn mới và hơi cồn sát khuẩn còn vương vất sau buổi sáng sắp đặt thiết bị.

``Gâu!''

Một tiếng sủa lanh lảnh, khỏe khoắn vang lên cắt ngang tiếng rù rì của chiếc quạt thông gió.

Từ lối mòn rải sỏi dẫn từ xưởng cơ khí sang phòng thí nghiệm, một bóng trắng lao tới thoăn thoắt. Đó là Tô, cô chó cái tròn bốn tuổi, đứa con đầu lòng của Mi. Khác với người mẹ nhỏ thó chân ngắn, Tô lại sở hữu vóc dáng dong dỏng cao, bộ khung xương săn chắc và bốn chiếc chân nhanh nhẹn của một chú chó cỏ bản địa, là kết quả của mối tình vụng trộm năm nào giữa Mi và chú chó cỏ nhà hàng xóm. Duy chỉ có bộ lông trắng muốt như bông gòn và đôi tai tam giác vểnh cao là giữ trọn vẹn nét duyên dáng từ mẹ. Cái đuôi xù lông cong vút như chiếc móc câu liên tục lắc lư sang hai bên, đôi mắt hạt nhãn đen láy sáng rực khi phát hiện bóng người quen thuộc trước bậc thềm cửa.

Ngay phía sau Tô chừng dăm sáu bước chân, một chú cún khác nhỏ nhắn hơn hẳn, vóc dáng mảnh dẻ với bốn chiếc chân ngắn củn và bộ lông mềm mượt như tơ lụa trắng pha vài vệt kem nhạt quanh sống mũi, đang thong thả bước từng nhịp khoan thai. Khác hẳn với vẻ hùng hổ, lăn xả của cô chị cùng mẹ khác lứa, cô bé cún này chỉ khẽ nghiêng đầu sang một bên, đôi mắt to tròn phẳng lặng như mặt hồ thu, đưa ánh nhìn chăm chú dò xét từng nhành cây ngọn cỏ xung quanh.

Đó là Ki, cô con gái út vừa bước sang tuổi thứ hai vào mùa xuân vừa rồi, mang đậm nét vóc dáng Bắc Kinh lai Nhật của mẹ Mi.

``Tô! Ki! Không được chạy lung tung vào khu để hóa chất đâu đấy nhé!''

Giọng nói trong trẻo, ngọt ngào của Suzuno vang lên từ phía bàn thao tác. Cô bé đang khoác chiếc áo blouse trắng tinh khôi dài quá gối, hai tay đeo găng cao su mỏng, cẩn thận đặt từng ống nghiệm thủy tinh vào giá đỡ gỗ sồi. Nghe tiếng móng chân cào sàn lẹp kẹp, Suzuno buông nhẹ chiếc kẹp gắp, mỉm cười xoay người lại.

Tô chẳng buồn nghe lời nhắc nhở mang tính học thuật ấy. Cô nàng cún trượt một vệt dài trên nền gạch men bóng loáng, đâm sầm đầu vào ống quần xắn cao của Suzuno, vừa rên ư ử nũng nịu vừa liên tục dùng chiếc mũi ươn ướt hếch nhẹ vào mu bàn tay cô đòi vuốt ve.

Trong khi đó, Ki chỉ lẳng lặng tiến lại gần, ngồi bệt xuống cách chân Suzuno nửa bước, cái đuôi khẽ quét nhẹ trên mặt sàn hai nhịp rồi ngước lên nhìn cô bằng ánh mắt thông thái đến lạ kỳ, hệt như một học giả nhí đang tò mò quan sát một đề tài khoa học thú vị.

``Hai đứa lại trốn bố chạy sang đây à?''

Một giọng nói trầm ấm quen thuộc cất lên từ ngoài hiên. Kuon bước vào, tay xách chiếc xô nhựa đựng giẻ lau và bộ tua-vít. Cậu vừa giúp Suzuno siết lại các chốt thép của khung cửa sổ chịu lực. Cậu Tổng quản đưa mắt nhìn hai cái đầu lông trắng đang vây quanh chân em gái, khóe môi khẽ cong lên một nét cười dịu dàng:

``Anh vừa thấy bố vác xẻng ở góc vườn, bảo sao bỗng dưng không thấy hai tên lính gác cổng đâu, hóa ra là mò sang hưởng điều hòa của viện nghiên cứu bên em rồi.''

Suzuno ngồi xổm xuống, cởi chiếc găng cao su đặt lên khay khử trùng, rồi đưa cả hai bàn tay dịu dàng xoa nhẹ hai bên mang tai xù bông của Tô. Cô cún sướng rơn, ngửa hẳn cổ ra sau, hai mắt nhắm tịt lại, bốn chân co giật nhè nhẹ vì khoái chí.

``Hôm nay trời oi bức quá mà, Onii-sama. Xưởng của bố máy cắt chạy suốt ngày nên bụi với nóng lắm, hai đứa lại nhớ hơi người, cứ quanh quẩn bên chân em mãi thôi. Em mang chúng sang đây vừa để có bạn trò chuyện lúc pha dung môi, vừa nhờ Tô với Ki kiêm luôn nhiệm vụ gác cửa cho phòng thí nghiệm.''

Suzuno vừa nói vừa đưa tay sang xoa đầu Ki. Cô bé cún nhỏ khẽ dụi nhẹ vầng trán mềm mại vào lòng bàn tay Suzuno, không hề vồ vập nhưng sự âu yếm lại bền bỉ và sâu lắng khôn cùng.

Ngồi trên bậc thềm cửa hiên nhìn bóng nắng xiên qua tán cây bồ đề cổ thụ, Kuon khẽ thở dài một hơi nhẹ bẫng, ánh mắt cậu bỗng trôi dạt về một góc vườn sau xưởng cơ khí phủ đầy bóng râm mát rượi cách đó vài chục mét:

``Nhanh thật đấy\dots\ Mới đó mà đã tròn hai tháng rồi.''

Nụ cười trên môi Suzuno khẽ chùng xuống. Đôi bàn tay đang vuốt ve lông Tô thoáng khựng lại một nhịp, đáy mắt ánh lên một vệt buồn thăm thẳm:

``Vâng\dots\ Đầu tháng Sáu, đúng ngày cơn mưa rào chuyển mùa đầu tiên trút xuống.''

Đầu tháng Sáu vừa qua, khi những cánh hoa anh đào cuối cùng của mùa xuân đã rụng sạch và những cơn mưa rào báo hiệu mùa hạ oi ả bắt đầu đổ về Kawachi, Mi - cô chó mẹ của Tô và Ki, người cận vệ thầm lặng của gia đình Hikawa suốt mười bốn năm trời - đã thanh thản trút hơi thở cuối cùng vì tuổi già.

Mười bốn năm đối với một đời người có thể chỉ là một đoạn thanh xuân, nhưng đối với một chú cún, đó là cả một kiếp luân hồi dài dằng dặc, trọn vẹn và đủ đầy.

Suzuno khẽ vuốt ve sống lưng ấm áp của Ki, giọng nhỏ nhẹ như tiếng gió thoảng:

``Từ bé em đã rất thích chó, yêu cái cảm giác được ôm một sinh linh trung thành vào lòng. Hồi mới được anh đón về căn nhà này, nhìn thấy xưởng nhà mình nuôi không chỉ một mà là tận ba con chó, em đã mừng rỡ biết bao nhiêu\dots\ Của đáng tội là khi em gặp Mi, con bé đã quá già rồi. Giống Bắc Kinh lai Nhật vốn dĩ vóc người đã nhỏ nhắn, đến lúc tuổi già mắt lại phủ lớp màng đục mờ, bốn chân run rẩy nhưng sớm nào con bé cũng cố gượng dậy lết ra bậc cửa xưởng nằm canh cho bố. Nhìn con bé kiên cường như thế, em vừa thương vừa cảm phục vô cùng.''

Kuon ngồi xuống bậc thềm gỗ bên cạnh em gái, duỗi dài hai chân, ánh mắt xa xăm nhớ lại:

``Bởi vì con bé không xem mình là một con thú cưng nuôi để giải trí. Từ ngày đầu tiên bước chân vào cổng xưởng cơ khí này, Mi đã tự xem mình là người nhà Hikawa rồi.''

Cậu Tổng khẽ bật cười, một nụ cười ấm áp xen lẫn chút bồi hồi khi những thước phim ký ức của mười bốn năm về trước bỗng chốc cuộn trào sống động trong tâm trí:

``Cái ngày đón con bé về\dots\ đúng là một chuyện cười ra nước mắt mà có lẽ đến già anh cũng không quên được.''

\ct

\begin{center}
    \uline{Mười bốn năm trước.}
\end{center}

Đó là thời điểm gia đình Hikawa vừa quyết định rời bỏ vùng thị trấn ven biển đầy nắng gió và vị mặn mòi của những con sóng để chuyển lên thủ đô Kawachi lập nghiệp được tròn hai năm. Ryujin dốc toàn bộ vốn liếng tích cóp sau những năm tháng bôn ba để cải tạo khu xưởng cơ khí ở vùng ngoại ô được truyền từ ông nội.

Năm ấy, Kuon mười hai tuổi, đang học cấp hai. Còn Ryuuhou khi ấy mới vừa tròn hai tuổi, chập chững biết đi những bước chân ngắn ngủn xiêu vẹo và cái miệng nhỏ chúm chím suốt ngày líu lo bập bẹ những câu từ ngọng nghịu.

Khu xưởng khi mới tiếp quản rộng thênh thang, nhiều ngóc ngách và đêm đêm gió lùa hun hút. Máy móc, sắt thép chất thành từng dãy cao ngất ngưởng, thỉnh thoảng lại có chuột cống và mèo hoang chạy rông làm rơi vỡ đồ đạc. Ông đã bàn với vợ mình: cần phải nuôi một chú chó để gác xưởng, vừa để canh giữ vật tư, vừa để đêm hôm sủa báo động nếu có kẻ gian nhòm ngó.

Vào một buổi sáng Chủ nhật trời thu trong vắt, ông cưỡi chiếc xe máy cũ chở vợ, còn Kuon thì dắt tay cô nhóc Ryuuhou hai tuổi đi bộ tới một tiệm bán thú cưng và chó giống nằm ở cuối con phố buôn bán sầm uất của thị trấn Kawachi.

Trong tâm tưởng của người chủ xưởng cơ khí - một người đàn ông miền Bắc dạn dày sương gió, từng trải qua lao động nặng nhọc - một chú chó giữ xưởng đạt chuẩn bắt buộc phải là những giống chó to lớn, hung dữ và oai vệ: hoặc là một chú Akita lông vàng vạm vỡ, hoặc là giống chăn cừu Đức lai với bốn chân to bản và tiếng sủa ồm ồm đủ khiến trộm vía kinh hồn bạt vía. Ông xăm xăm rảo bước dọc theo những chiếc lồng sắt lớn ở dãy ngoài cùng, liên tục ngắm nghía, huýt sáo thử lòng dũng cảm của mấy chú chó giống to lớn.

Thế nhưng, số phận của gia đình Hikawa lại rẽ sang một hướng hoàn toàn trớ trêu, chỉ vì sự bướng bỉnh ngây thơ của một cô bé hai tuổi.

Trong khi người lớn đang mải mê bàn luận về giống chó săn, Ryuuhou nhỏ thó, mặc chiếc yếm bò màu xanh dương phối áo thun vàng, hai bím tóc ngắn ngủn buộc vểnh sang hai bên như hai cái sừng dê, đã tự mình lạch bạch chạy tót vào khu vực chuồng kính dành cho các giống chó cảnh nhỏ con.

Đó là một chiếc lồng kính đặt ở tầng thấp nhất, cách mặt đất chưa đầy ba tấc.

Bên trong lồng, những chú chó con thuộc các giống Poodle, Phốc sóc hay Chihuahua đều đang cuộn tròn ngủ say hoặc sủa nhặng xị vì ồn ào. Duy chỉ có một sinh linh nhỏ bé đang tỉnh táo.

Đó là một cô cún lai đặc biệt: nửa mang dòng máu chó Bắc Kinh (Pekingese) quý phái với khuôn mặt ngắn tũn, vầng trán dô tròn và chiếc mũi hếch ngộ nghĩnh; nửa mang dòng máu chó Nhật Bản (Japanese Chin) với bộ lông trắng muốt như tuyết đầu mùa, điểm xuyết những đốm màu kem thanh nhã quanh hai vành tai cụp mềm và chiếc đuôi uốn cong xòe lông như cánh quạt lông công.

Khoảnh khắc cô cún con nhìn thấy Ryuuhou đứng ngơ ngác bên ngoài tấm kính, đôi mắt đen láy to tròn, ươn ướt của nó bỗng sáng rực lên.

Không một chút sợ sệt, cô cún lon ton chạy thẳng về phía góc lồng nơi Ryuuhou đang đứng. Nó đứng thẳng bằng hai chân sau, hai chiếc chân trước ngắn tũn cào cào nhẹ nhàng lên mặt kính trong suốt, cái đầu tròn xoe nghiêng hẳn sang một bên, cái miệng nhỏ hé mở để lộ đầu lưỡi hồng hào khẽ liếm láp như muốn chạm vào những ngón tay múp míp của cô bé ngoài lớp kính.

Ryuuhou đứng sững lại.

Đôi mắt to tròn ngây thơ của cô bé hai tuổi nhìn chăm chú vào cô cún nhỏ bên trong. Rồi bỗng nhiên, con bé toét miệng cười giòn tan, vỗ hai tay vào nhau bôm bốp: ``Cún! Cún trắng!''

Nghe tiếng em gái cười nắc nẻ, Kuon vội buông tay mẹ chạy lại: ``Ryuuhou, không được đập vào kính của người ta, kẻo nó giật mình đấy!''

Thế nhưng, vừa thấy Kuon tiến lại gần, cô cún nhỏ bên trong chẳng những không sợ mà còn quẫy chiếc đuôi tít mù đến mức cả phần hông tròn lẳn rung rinh theo. Nó cất lên một tiếng sủa non nớt: ``Ăng!'' - một âm thanh thanh mảnh, êm tai như tiếng chuông gió rung nhẹ trong chiều thu.

Đúng lúc đó, Ryujin bước tới, tay đang cầm tờ giấy giới thiệu giống chó săn Rottweiler: ``Sao thế mấy đứa? Đi ra đằng kia xem mấy con béc-giê lai đi, to khỏe lắm, sau này trông kho mới yên tâm được.''

Nào ngờ, vừa nghe bố nói thế, Ryuuhou lập tức ôm chặt lấy mép kính của chiếc chuồng nhỏ, hai chân dính chặt xuống sàn gạch như đóng đinh. Con bé phồng má, chu môi, hai mắt ngân ngấn nước nhưng nét mặt thì cương quyết đến mức vô lý: ``Hông! Nuôi cún này cơ! Cún này xinh! Bế cún này về nhà!''

``Kìa Ryuuhou, con này là chó cảnh, bé như nắm xôi thế này thì trông nhà trông xưởng sao được con?'' Ryujin dở khóc dở cười xoa đầu con gái, cố gắng giải thích bằng lý lẽ người lớn. ``Nuôi chó là để giữ xưởng cho bố, con này ra đường khéo mèo hoang nó tát cho chạy mất dép ấy chứ!''

Nhưng với trí tuệ non nớt của một đứa trẻ hai tuổi, lý lẽ logic của người lớn hoàn toàn là con số không tròn trĩnh.

Ryuuhou bắt đầu dẩu môi, nước mắt cá sấu lăn dài trên đôi má phúng phính. Con bé ngồi bệt xuống sàn đá hoa, hai tay bám riết lấy khung kính không chịu buông, vừa mếu máo vừa gào lên bài ca muôn thuở của trẻ con:

``Hông chịu đâu! Bố bắt cún to con sợ! Bế cún trắng! Cún trắng liếm tay Ryuuhou rồi! Cún trắng là bạn của Ryuuhou!''

Bên trong lồng kính, dường như cũng cảm nhận được ``đồng minh'' nhỏ tuổi đang khóc vì mình, cô cún Bắc Kinh lai Nhật cũng sủa nhặng xị lên vài tiếng đầy đồng cảm, rồi áp sát mõm vào khe hở thông khí của lồng, cố vươn cái lưỡi nhỏ xíu ra ngoài như muốn lau nước mắt cho cô bé hai tuổi.

Trước cảnh tượng dở khóc dở cười ấy, Kaien bật cười chịu thua, huých nhẹ khuỷu tay vào sườn ông chồng to xác: ``Kìa anh! Con bé có duyên với nó đấy. Người ta bảo chó chọn chủ chứ chủ mấy khi chọn được chó. Trông con cún mắt sáng, lanh lợi thế kia, tuy nhỏ con nhưng có linh tính đấy.''

Kuon khi ấy cũng nhìn cô cún nhỏ với ánh mắt đầy thiện cảm. Cậu ngẩng lên nhìn bố, nói thêm vào: ``Bố ơi, xưởng mình có khung rào kín rồi mà. Nuôi con cún này vừa để em Ryuuhou có bạn chơi ở nhà, vừa tập cho nó biết giữ xưởng từ bé cũng được mà bố.''

Bị cả vợ lẫn hai đứa con ``tổng tấn công'', người đàn ông trụ cột của gia đình Hikawa chỉ biết đưa tay gãi đầu gãi tai, thở dài một hơi bất lực nhưng khóe mắt lại tràn đầy vẻ cưng chiều: ``Thôi được rồi! Chiều cô út hết đấy nhé! Sau này nó mà không biết sủa trộm thì cô út liệu hồn đi dọn phân cho nó nghe chưa!''

Ryuuhou nghe thấy thế thì lập tức nín bặt, nước mắt nước mũi còn tèm lem trên mặt nhưng miệng đã cười toe toét, nhảy cẫng lên ôm chầm lấy chân bố.

Khi nhân viên cửa hàng mở lồng kính bế cô cún ra, trao vào vòng tay vụng về của cô bé hai tuổi, cô cún chẳng hề giãy giụa hay sợ hãi. Nó ngoan ngoãn thu mình trong vòng tay bé xíu của Ryuuhou, đưa chiếc mũi ươn ướt hít hà khắp cổ cô bé rồi khẽ dụi đầu vào hõm vai con bé, phát ra những tiếng gừ gừ êm ái như một động cơ mini.

Trên đường về nhà, ngồi gọn trong lòng Kuon ở phía sau xe máy, con cún nhỏ liên tục liếm ngón tay cậu và và cô em gái.

``Thế giờ đặt tên nó là gì đây?'' Ryujin vừa lái xe vừa hỏi với ra phía sau.

Ryuuhou đang ngậm chiếc kẹo mút dở, nhìn cái mặt ngắn tũn đầy ngộ nghĩnh của cô cún lông trắng, liền bập bẹ thốt lên từ ngữ ngẫu nhiên đầu tiên nảy ra trong đầu non nớt: ``Mi! Mi mi!''

``Mi à? Nghe cụt lủn thế?'' Kuon nhướng mày trêu em.

``Mi! Là Mi-chan!'' Ryuuhou dẩu môi khăng khăng, hai tay giữ chặt lấy tai cún con như để khẳng định chủ quyền tối cao.

Và thế là, cô chó giống Bắc Kinh lai Nhật bé nhỏ ấy chính thức gia nhập hộ khẩu nhà Hikawa với cái tên giản dị: Mi.

\ct

Thế nhưng, chuỗi ngày rộn rã tiếng cười của gia đình Hikawa khi đón thành viên bốn chân mới về nhà lại chẳng kéo dài được bao lâu.

Ngay vào cuối năm 2008 ấy, một cơn bão ngầm bất ngờ ập xuống mái ấm nhỏ. Mâu thuẫn sâu sắc giữa hai dòng họ - nhà Hikawa bên nội và nhà Musada bên ngoại - bùng phát dữ dội quanh những bất đồng gay gắt về công việc làm ăn và quan điểm sống. Những cuộc cãi vã căng thẳng giữa người lớn cứ thế nối tiếp nhau, vượt quá tầm kiểm soát và sức chịu đựng của hai bên.

Cái kết đắng cay không ai mong muốn đã xảy ra: bố Ryujin và mẹ Kaien quyết định tạm thời ly thân.

Mẹ Kaien ôm đứa con gái nhỏ Ryuuhou rời khỏi Kawachi, trở về nương náu ở nhà ông bà ngoại để tiện bề chăm sóc con thơ. Kuon, đứa con trai cả mười hai tuổi, ban đầu cũng theo mẹ và em gái về bên ngoại. Thế nhưng, chỉ sau vài tuần ngắn ngủi sống trong bầu không khí ngột ngạt của sự chia cắt, hình ảnh người cha lầm lũi một thân một mình giữa gian xưởng cơ khí lạnh ngắt, vắng tanh không một bóng người cứ ám ảnh tâm trí cậu bé mười hai tuổi. Bố lặn lội sang đón Kuon, và cậu đã quyết định theo bố trở về lại xưởng cơ khí Kawachi để hai cha con nương tựa vào nhau.

Suốt gần ba năm đằng đẵng từ cuối năm 2008 đến năm 2011, căn xưởng cơ khí rộng thênh thang chỉ còn lại hai chiếc bóng đơn côi của hai cha con.

Mẹ và em gái vắng nhà, tiếng bi bô cười đùa của cô nhóc Ryuuhou tắt lịm, căn nhà nhỏ ở Anwa vào những đêm đông lạnh buốt đến thấu xương. Giữa chuỗi ngày u ám và trống trải tột cùng ấy, cô cún nhỏ Mi vô tình trở thành sợi dây níu giữ chút ấm áp gia đình còn sót lại.

Bố bận tối tăm mặt mũi với việc vực dậy xưởng cơ khí sau biến cố gia đình, ngày đêm miệt mài bên máy cưa, máy cắt để kiếm tiền trang trải nợ nần và gửi tiền chu cấp sang bên ngoại cho Ryuuhou. Mọi gánh nặng chăm sóc Mi gần như dồn hết lên đôi vai gầy của cậu bé Kuon từ năm mười hai đến mười lăm tuổi.

Vốn là người từng lớn lên ở làng quê và có bề dày kinh nghiệm nuôi dạy chó săn, người bố đã kiên nhẫn chỉ bảo cho cậu từng chút một: từ cách nấu nồi cháo tấm trộn đầu cá băm nhuyễn sao cho đủ chất mà không hóc xương, cách pha nước ấm tắm cho Mi giữa mùa đông rồi sấy thật khô bộ lông dài để tránh nấm da, cách nhổ ve rận ở kẽ tai, cho tới việc dắt Mi đi dạo mỗi chiều quanh các ngõ hẻm để rèn luyện thể lực.

Kuon học rất nhanh. Dưới sự dạy dỗ nghiêm khắc mà ân cần của bố, cậu Tổng quản nhí ngày ngày tự tay làm sạch chuồng, bưng từng bát cơm, chải từng lọn lông rối cho Mi. Cậu còn kiên nhẫn huấn luyện cho cô cún từng khẩu lệnh cơ bản: từ việc ngồi yên, nằm chờ, bắt tay, cho tới việc nhận biết tiếng bước chân người lạ và sủa báo động mỗi khi có kẻ bén mảng lại gần kho đá lát.

Những buổi chiều muộn sau giờ tan học, Kuon lại ngồi bên chiếc bàn gỗ cũ ở góc xưởng cặm cụi làm bài tập, còn Mi thì ngoan ngoãn nằm gác mõm lên mu bàn chân cậu, dùng thân nhiệt ấm áp sưởi ấm cho cậu qua những mùa đông buốt giá. Mỗi khi Kuon buồn tủi vì nhớ mẹ và nhớ em gái, Mi lại khẽ khàng vươn đầu lưỡi liếm nhẹ lên những giọt nước mắt lăn trên má cậu, đôi mắt nâu hạt nhãn ngước nhìn đầy thấu hiểu như muốn nói: ``Đừng khóc, có em ở đây với anh rồi.''

Chính trong ba năm gian khó ấy, mối dây liên kết giữa Kuon và Mi đã vượt xa khỏi tình cảm giữa chủ và vật nuôi thông thường. Con bé trở thành người bạn tri kỷ, người lính gác trung thành và là nhân chứng sống chứng kiến những nốt thăng trầm lặng lẽ nhất của hai cha con nhà Hikawa.

Phải mãi đến năm 2011 - khi Ryuuhou tròn năm tuổi - sau bao nhiêu nỗ lực hòa giải của họ hàng và sự kiên trì hàn gắn của Ryujin, mâu thuẫn giữa hai gia đình nội ngoại mới chính thức được tháo gỡ êm thấm. Mẹ Kaien dắt tay cô bé Ryuuhou năm tuổi trở về lại mái ấm Kawachi trong sự vỡ òa hạnh phúc của cả gia đình.

Ngày mẹ và em gái bước chân qua cổng xưởng, Mi - khi ấy đã là một cô chó ba tuổi trưởng thành, bộ lông trắng muốt xù bông óng ả và đôi mắt sáng quắc uy nghi - lập tức lao ra đón từ cửa nhà. Dù đã xa cách gần ba năm trời, Mi vẫn nhận ra ngay mùi hương của cô chủ nhỏ từng nằng nặc khóc đòi bế mình năm nào. Cô cún mừng rỡ đến mức nhảy cẫng lên quẫy đuôi tít mù, liếm láp khắp bàn tay nhỏ nhắn của Ryuuhou năm tuổi, rộn rã cất tiếng sủa chào mừng người nhà trở về đoàn tụ.

\begin{center}
    \uline{Kết thúc hồi tưởng.}
\end{center}

Gió chiều thổi lộng qua khung cửa sổ phòng lab, làm lay động những vạt nắng vàng óng ả đang trải dài trên mặt bàn đá hoa cương.

Nghe Kuon kể lại quãng thời gian ly thân đầy thăng trầm thuở nào, Suzuno không kìm được xúc động, khẽ siết nhẹ lấy bàn tay anh trai: ``Thảo nào\dots\ Mỗi lần Ryuuhou khoe về Mi, mặt em ấy vênh lên bao nhiêu thì ánh mắt anh Kuon nhìn Mi lại sâu lắng bấy nhiêu. Hóa ra trong ba năm xa vắng ấy, Mi chính là điểm tựa tinh thần lớn nhất của hai bố con anh.''

``Ừ,'' Kuon cười nhẹ, ngón tay cậu khẽ gãi gãi cằm Tô, ``con bé đã cùng anh và bố đi qua những tháng ngày tăm tối nhất của gia đình. Và đến khi trưởng thành, Mi lại mang đến cho xưởng những niềm vui bất ngờ khác.''

Cậu Tổng đưa mắt nhìn sang Tô - cô chó đang lim dim ngủ với vóc dáng cao ráo của loài chó cỏ bản địa: ``Năm Mi mười tuổi - cái tuổi mà đối với loài chó cảnh đã bước sang sườn dốc bên kia của cuộc đời - đùng một cái, con bé lại dính vào một mối tình vụng trộm đầy trớ trêu. Chẳng biết bằng cách nào, cô nàng tiểu thư Bắc Kinh chân ngắn nhà mình lại trèo rào sang `tâm sự' với chú chó cỏ giống đực vạm vỡ của xưởng cơ khí bên cạnh. Thế là kết quả, con bé vượt cạn thành công và sinh ra cô nàng Tô này đây.''

Suzuno bật cười khúc khích, đưa tay xoa nhẹ bờ vai săn chắc của Tô: ``Bảo sao Tô lại có bốn chân dài rỏng rảnh, mõm thon dài và dáng chạy khỏe khoắn hệt như chó cỏ như thế này! Nhưng mà công nhận nhờ lai giống chó cỏ mà Tô khỏe mạnh, ít ốm đau hơn hẳn mẹ Mi anh nhỉ.''

``Khỏe quá mức ấy chứ,'' Kuon cười lắc đầu, ``Tô sinh ra thừa hưởng trọn vẹn sự dẻo dai của bố và nét đẹp lông trắng của mẹ, tính nết thì nghịch ngợm, xốc nổi y hệt như con bé Ryuuhou vậy. Mi chăm Tô chu đáo từng ly từng tí, dạy con bé từng thế đứng, từng tiếng sủa báo trộm trong xưởng.''

Kuon chuyển ánh mắt sang Ki - cô cún nhỏ đang ngoan ngoãn nằm cuộn tròn bên chân Suzuno: ``Rồi hai năm sau đó, khi Mi đã tròn mười hai tuổi - độ tuổi xưa nay cực hiếm đối với loài chó - con bé lại sinh tiếp lứa thứ hai, chính là Ki. Lần này, Ki lại thừa hưởng trọn vẹn nét quý phái của mẹ: chân ngắn, mặt ngắn tũn ngộ nghĩnh, tính tình thì trầm tĩnh, biết quan sát và dịu dàng hệt như mẹ Mi thời trẻ.''

Ba mẹ con nhà cún cứ thế quấn quýt bên nhau trong căn xưởng cơ khí rộn rã tiếng cười, cùng nhau gác cổng, cùng nhau tắm nắng dưới gốc cây rẻ quạt nơi góc vườn xưởng.

Thế nhưng, quy luật khắc nghiệt của thời gian chưa từng bỏ quên bất kỳ một sinh mệnh nào.

Bước sang tuổi mười bốn, đôi mắt của Mi mờ dần bởi căn bệnh đục thủy tinh thể tuổi già, lớp lông trắng bóng mượt ngày nào dần ngả sang màu tro xám và xơ xác. Những bước chân thoăn thoắt năm xưa trở nên chậm chạp, run rẩy trên nền sỏi đá lởm chởm. Mi không còn đủ sức phóng qua bậc cửa tam cấp để đón Kuon mỗi chiều đi làm về nữa.

Thế nhưng, sự trung thành và thói quen mười bốn năm dường như đã ăn sâu vào tủy tủy. Sáng nào cũng vậy, dù xương khớp đau nhức vì thời tiết chuyển mùa, cô chó già vẫn kiên trì lết từng bước nặng nhọc từ góc chuồng ra nằm phục dưới tán cây rẻ quạt, nơi có tầm nhìn thẳng ra cổng xưởng và cửa phòng thí nghiệm đang xây dở của Suzuno.

Và rồi, buổi sáng định mệnh đầu tháng Sáu ấy đã đến.

Đó là một ngày nắng hanh hao sau cơn mưa rào chuyển mùa rả rích suốt đêm. Suzuno mang bát cháo thịt băm thơm phức nấu riêng cho Mi ra vườn, nhưng cô chó già không còn cất tiếng sủa khẽ khàng đón chào như mọi khi nữa.

Mi nằm cuộn tròn bình yên trên thảm cỏ xanh mướt dưới gốc cây rẻ quạt cổ thụ, đầu tựa nhẹ lên hai chân trước, đôi mắt nhắm nghiền thanh thản như đang chìm vào một giấc mộng dài tuyệt đẹp sau mười bốn năm canh gác không một ngày ngơi nghỉ.

Bên cạnh mẹ, Tô và Ki nằm nép sát vào lưng Mi, liên tục dùng mũi hếch nhẹ vào tai mẹ, phát ra những tiếng rên ư ử nghẹn ngào, đau xót như muốn đánh thức người mẹ già dậy ăn bữa sáng.

Nhưng Mi đã vĩnh viễn không bao giờ tỉnh lại nữa.

Một giọt nước mắt trong veo khẽ lăn dài trên gò má trắng ngần của Suzuno, rơi khẽ xuống bộ lông tơ mềm mại của Ki.

Cô bé cún dường như cảm nhận được nỗi buồn trong lòng người chị cả, khẽ vươn đầu lưỡi nhỏ liếm nhẹ lên mu bàn tay Suzuno, đôi mắt hạt nhãn ánh lên vẻ an ủi sâu sắc.

``Bố đã đích thân dùng chiếc xẻng quen thuộc chôn cất Mi ngay dưới gốc cây rẻ quạt ấy\dots'' Suzuno khẽ thì thầm, bàn tay siết nhẹ lấy tấm lưng nhỏ của Ki. ``Bố không cho ai đụng vào. Bố bảo, Mi đã nằm gác ở cái bóng râm ấy suốt mười bốn năm, thì hãy để con bé được yên nghỉ mãi mãi ở nơi nó thân thuộc nhất. Bố còn đặt một tảng đá hoa cương trắng khắc tên Mi lên làm mộ chí nữa.''

Kuon gật đầu, đưa tay vỗ nhẹ bờ vai em gái: ``Hôm ấy cả nhà đều có mặt. Cả Kaiho vốn lạnh lùng ít nói là thế, mà lúc nhìn thấy tảng đá hoa cương được đặt xuống cũng phải quay mặt đi chỗ khác giấu giọt nước mắt. Còn con bé Ryuuhou thì khóc ngất đi suốt cả buổi chiều, cứ ôm khư khư chiếc vòng cổ có gắn chuông đồng cũ kỹ của Mi không chịu buông.''

Sau khi Mi qua đời, căn xưởng cơ khí bỗng trở nên trống trải đến lạ kỳ. Dù vẫn còn tiếng búa gõ, tiếng máy cắt gọt rầm rập của Ryujin và thợ phụ, nhưng khoảng sân xưởng dường như đã đánh mất đi một phần linh hồn ấm áp nhất.

Tô và Ki mồ côi mẹ, suốt tuần lễ đầu tiên cứ nằm lỳ bên tảng đá hoa cương trắng dưới gốc cây, chẳng buồn ăn uống, ánh mắt cứ ngơ ngác nhìn ra cổng sắt như chờ đợi một bóng hình thân thuộc quay về.

Chính vì xót xa trước nỗi đau của hai chú cún nhỏ, Suzuno và Naomi đã xin phép bố được chăm sóc hai đứa. Hằng ngày, sau những giờ học và nghiên cứu, Suzuno lại cùng Naomi dắt Tô và Ki đi dạo quanh khu đất hoang, đưa chúng sang phòng thí nghiệm vừa xây xong để bầu bạn, chăm chút từng bữa ăn dinh dưỡng và chải chuốt bộ lông trắng muốt cho chúng.

Dần dà, sự ấm áp của gia đình Hikawa đã xoa dịu đi nỗi mất mát trong lòng hai chú chó nhỏ.

Tô đã lấy lại được vẻ hiếu động, tinh nghịch ngày nào, tự nguyện biến mình thành ``đội trưởng đội cảnh vệ'' chạy lăng xăng khắp xưởng và phòng thí nghiệm; còn Ki thì trở thành cái đuôi nhỏ dịu dàng của Suzuno, lặng lẽ ngồi bên bàn thí nghiệm canh chừng từng ống nghiệm hóa chất cho cô.

Suzuno hít một hơi thật sâu, lau đi giọt nước mắt còn vương trên khóe mi, nhìn hai chú cún đang quấn quýt dưới chân mình với ánh mắt tràn đầy sự yêu thương: ``Anh Kuon này\dots\ Em tin rằng linh hồn của Mi chưa bao giờ rời xa gia đình mình đâu. Bằng một cách nào đó, bà mẹ phi thường ấy vẫn đang dõi theo chúng ta, che chở cho căn xưởng của bố, và gửi gắm trọn vẹn tình yêu thương cùng sự trung thành của mình vào hai đứa con nhỏ này.''

Kuon mỉm cười, ánh mắt cậu hướng ra khoảng vườn ngập nắng ngoài hiên, nơi tảng đá hoa cương trắng dưới gốc cây rẻ quạt vẫn đang tỏa sáng lấp lánh dưới những tia nắng vàng ấm áp: ``Ừ. Gia đình Hikawa chúng ta\dots\ chưa từng bỏ rơi bất kỳ ai, và cũng sẽ không bao giờ quên đi những người đã hết lòng vì mái ấm này.''

Cậu Tổng khẽ cúi người, vỗ nhẹ vào đầu Tô: ``Nào, hai cô lính gác nhỏ! Nghỉ giải lao thế là đủ rồi. Giờ thì phụ chị Suzuno canh chừng cửa lab cho cẩn thận nhé, chiều nay anh sẽ bảo bà Kaien thưởng cho mỗi đứa một khúc xương ống hầm thật to!''

``Gâu! Gâu!''

Tô sủa vang hai tiếng đầy tự hào như một lời hứa danh dự của người lính, rồi cùng Ki nhảy cẫng lên, rộn rã quấn quýt quanh bóng dáng của hai anh em dưới bầu trời mùa hạ thanh bình của thủ đô Kawachi.

\ct

Khoảng hơn bốn giờ chiều, khi cái nắng gay gắt bắt đầu dịu dần, nhường chỗ cho những ngọn gió nam mát rượi thổi qua bãi đất hoang, Kuon chào hai cô em gái rồi quay trở lại gian xưởng cưa để phụ bố kiểm kê đợt trục thép vừa nhập về.

Trước cửa phòng thí nghiệm, cánh cửa lưới chống côn trùng khẽ kêu cọt kẹt. Naomi - cô bé được Kuon và Suzuno đưa xuống chơi - bước ra từ lối đi nhỏ nối với ký túc xá, hai tay cẩn thận bê một chiếc khay gỗ nhỏ. Trên khay là đĩa dưa hấu đỏ au đã được cắt thành từng miếng tam giác ướp lạnh ngọt lịm, cùng một chiếc hộp gỗ mài nhẵn đựng dụng cụ vệ sinh cá nhân.

``Chị Suzuno\dots\ Em mang dưa hấu sang cho chị này,'' cô bé khẽ cất giọng, nhỏ nhẹ và êm ái như tiếng nước suối róc rách.

``A, cảm ơn Naomi nhé! Đúng lúc chị vừa pha chế xong lượt dung dịch đệm, khát khô cả cổ rồi đây.''

Suzuno tháo chiếc áo blouse treo lên giá ngoài hiên, đón lấy đĩa dưa hấu từ tay em gái, khẽ chia cho Naomi một miếng lớn. Mùi dưa hấu ngọt mát lập tức thu hút hai cái mũi thính nhạy của hai cô cún. Tô thoăn thoắt phi lại gần, hai chân trước đặt hẳn lên mép bậc thềm gỗ, cái đuôi quẫy tít mù, mắt nhìn chăm chú vào miếng dưa hấu trên tay Naomi đầy thèm thuồng.

Naomi mỉm cười hiếm hoi, bẻ một mẩu ruột đỏ không hạt đưa ra trước mũi Tô: ``Cho Tô này\dots\ Nhưng phải ăn từ từ thôi nhé, không được nuốt chửng đâu.''

Tô ngoan ngoãn liếm nhẹ lấy miếng dưa từ đầu ngón tay Naomi, nhai sột soạt đầy thích thú. Trong khi đó, Ki chỉ thong thả bước lại, ngồi nép bên tà váy của cô, đưa chiếc đầu nhỏ áp vào mu bàn chân cô bé như một lời chào quen thuộc. Naomi cúi người ôm lấy Ki, bàn tay nhỏ nhắn vuốt ve lớp lông mềm mịn của cô bé cún, mắt ánh lên vẻ thanh thản lạ kỳ. Giữa một đứa trẻ trầm lặng ít nói như Naomi và cô cún nhỏ Ki dường như luôn tồn tại một sợi dây liên kết vô hình, chẳng cần lời nói mà vẫn thấu hiểu trọn vẹn tâm ý của nhau.

Ăn xong miếng dưa mát lạnh, Suzuno kéo chiếc ghế đẩu bằng gỗ sồi ra khoảng sân hiên râm mát dưới giàn hoa giấy tím biếc, mở chiếc hộp gỗ mà cô em út vừa mang sang: ``Tranh thủ trời còn sáng, chị em mình chải lông cho hai đứa nó nhé. Mấy hôm nay giao mùa, lông tơ của Tô rụng nhiều quá, cứ bám đầy vào ống quần của bố mỗi khi ra xưởng.''

Trong chiếc hộp gỗ cũ kỹ là vài chiếc bấm móng, chai xịt dưỡng lông chiết xuất thảo mộc và một chiếc lược chải tóc cán gỗ mun đã ngả màu thời gian. Đó là chiếc lược chải dùng chung mà cả đại gia đình Hikawa vẫn thường tiện tay để ở chiếc bàn con ngay lối ra vào tầng một ký túc xá.

Từ anh cả Kuon mỗi sáng trước khi khoác áo đi làm, chị hai Kaiho vuốt lại mái tóc dài sau những giờ khổ luyện thể lực, cô nàng Ryuuhou ngổ ngáo chải vội mái tóc rối bời trước khi bật stream, cho tới Suzuno và Naomi ngày ngày chăm chút trước gương, chiếc lược gỗ ấy đã đi cùng nhịp sống thường nhật của năm anh chị em suốt bao tháng ngày. Trên chân những chiếc răng lược bằng thép bọc đầu tròn nhẵn bóng, dường như vẫn còn vương vất vài sợi tóc màu đen tuyền và nâu sẫm của các thành viên trong gia đình mà Suzuno chưa kịp gỡ hết.

``Tô, lại đây nào cô nương!'' Suzuno vỗ nhẹ lên đùi ra hiệu.

Cô nàng chó cỏ lai vốn hiếu động là thế, ban đầu còn ngúng nguẩy định chạy đuổi theo mấy con bướm vàng ngoài bãi cỏ, nhưng vừa thấy Suzuno cầm chiếc lược quen thuộc lên thì lập tức ngoan ngoãn tiến lại, ngoan ngoãn nằm bẹp bốn chân trên nền gạch men mát rượi.

Suzuno ngồi thụp xuống, tay trái giữ nhẹ vai Tô, tay phải đưa những đường lược nhịp nhàng, êm ái miết dọc theo chiều lông từ gáy xuống tận chóp đuôi. Chiếc lược kim loại đi tới đâu, những sợi lông chết và lớp lông tơ xù trắng muốt của Tô theo tới đó. Được gãi đúng chỗ ngứa dọc sống lưng, Tô sướng rơn người, hai mắt lim dim nhắm tịt lại, chiếc mõm dài thon khẽ thở ra những tiếng phì phò đầy mãn nguyện, cái đuôi thỉnh thoảng lại khẽ đập nhịp hai cái xuống sàn.

Từng búi lông tơ trắng muốt như tuyết được chiếc lược giữ lại, bện chặt lấy những sợi tóc người vương sẵn trên chân răng lược. Suzuno vừa chải vừa mỉm cười vuốt ve: ``Ngoan lắm\dots\ Lông Tô dày và khỏe thật đấy. Nhìn cái lưng săn chắc này xem, y hệt như mấy chú chó săn dũng cảm của các bộ lạc thảo nguyên vậy.''

Chải xong cho Tô, Suzuno vỗ nhẹ hông cô nàng: ``Xong rồi nhé, giờ đến lượt cô em gái Ki nào.''

Tô đứng dậy, lắc mạnh toàn thân từ đầu đến đuôi làm phát ra những tiếng phành phạch giòn giã, rồi lùi ra một góc nằm sưởi nốt chút nắng ấm cuối ngày.

Ki bước tới, thong thả tựa chiếc đầu ngắn tũn vào lòng bàn tay Suzuno. Khác với bộ lông dày và cứng cáp của Tô, lông của Ki dài và mềm như sợi tơ tằm, mang sắc trắng tinh khiết pha chút ánh kem thanh lịch của giống Bắc Kinh lai Nhật.

Suzuno đổi góc lược, nhẹ nhàng gỡ từng lọn lông rối quanh hai vành tai cụp và vùng yếm ngực trắng như tuyết của Ki. Naomi ngồi bên cạnh, nghiêng đầu ngắm nhìn chị gái thao tác, khẽ thầm thì: ``Ki ngoan thật đấy chị nhỉ\dots\ Ánh mắt của Ki lúc nào cũng tĩnh lặng như thể em ấy hiểu hết mọi chuyện xung quanh vậy.''

``Ừ,'' Suzuno dịu dàng đáp, tay vẫn thoăn thoắt đưa từng đường lược tỉ mỉ, ``Ki giống hệt như một quý cô nhỏ vậy, luôn điềm đạm và biết lắng nghe. Mỗi lần chị làm thí nghiệm thất bại, nhìn vào mắt Ki là lòng chị lại bình yên trở lại.''

Những sợi lông tơ mềm mại của Ki rụng ra, một lần nữa lại quấn chặt lấy các răng lược gỗ, hòa quyện không thể tách rời với những sợi tóc của anh em nhà Hikawa. Suzuno dùng đầu ngón tay gỡ toàn bộ búi lông trắng muốt dính tóc ấy ra khỏi chân lược, vo tròn lại thành một búi nhỏ bằng quả cau, đặt tạm lên chiếc đĩa sứ khử trùng trên bàn thao tác ngoài hiên:

``Chút nữa dọn lab xong, chị sẽ mang cuộn lông này vứt vào thùng rác hữu cơ, chứ để gió thổi bay lung tung lại dính vào quần áo.''

Cô bé hoàn toàn không thể ngờ rằng, búi lông tưởng chừng vô hại và tầm thường ấy - mang theo sự giao thoa kỳ diệu giữa tế bào của ba mẹ con nhà cún và các chuỗi mã ADN trứ danh của năm người con nhà Hikawa - lại chính là chiếc chìa khóa định mệnh sắp sửa làm đảo lộn số phận của cả gia đình.

\ct

Ráng chiều buông xuống ngoại ô Kawachi nhanh như một cái chớp mắt.

Mặt trời đỏ rực như một quả cầu lửa khổng lồ từ từ lặn khuất sau những rặng núi phía tây, dát một lớp ánh sáng màu cam đồng rực rỡ lên khắp mái tôn của xưởng cơ khí và khu phòng thí nghiệm phía sau. Gió hoàng hôn thổi qua, lay động những tán lá xanh rì của cây rẻ quạt cổ thụ, khiến những tia nắng muộn lấp lánh như dát bạc trên bề mặt tảng đá hoa cương trắng khắc tên Mi.

``Tô! Ki! Hai đứa lại đây nào!''

Suzuno đứng trước bậc thềm phòng thí nghiệm, tháo bỏ chiếc găng tay, mỉm cười cất tiếng gọi lớn ra góc vườn.

Dưới gốc cây rẻ quạt, nơi tảng đá hoa cương vẫn còn lưu giữ chút hơi ấm cuối cùng của ngày hè, hai bóng trắng đang nằm sóng đôi bên nhau bỗng giỏng tai lên. Tô bật dậy trước, bốn chân thoăn thoắt phi như bay qua thảm cỏ dại; Ki theo sau với những nhịp bước ngắn nhưng dứt khoát và nhanh nhẹn. Cả hai cô cún cùng chạy ùa về phía hai chị em Suzuno và Naomi, cái đuôi quẫy tít mù đầy rộn rã.

Từ phía cửa ngách thông sang xưởng cơ khí, tiếng còi báo hiệu hết ca làm việc vang lên giòn giã. Bố Ryujin trong bộ quần áo thợ bảo hộ còn lấm lem dầu máy, lau mồ hôi trên trán, cất giọng sang sảng vọng sang: ``Mấy đứa ơi! Dọn dẹp phòng nghiên cứu rồi chuẩn bị xe cộ về nhà ăn cơm thôi! Mẹ vừa gọi điện báo cơm sườn nướng với canh chua xong xuôi cả rồi đấy!''

``Vâng ạ! Chúng con sang ngay đây bố ơi!'' Suzuno đáp lại bằng chất giọng trong vắt.

Cô bé quay sang khóa cẩn thận cánh cửa kính của phòng thí nghiệm, tắt hệ thống điện ngoài hiên, cúi người xoa đầu dặn dò Tô và Ki ngoan ngoãn gác xưởng qua đêm, rồi cùng Naomi rảo bước qua lối mòn rải sỏi hướng về phía gian xưởng cơ khí. Ở đó, bố Ryujin và anh Kuon đã dắt sẵn hai chiếc xe máy ra cổng, nổ máy giòn tan chờ chở mấy đứa vượt qua bốn phường trong làn gió chiều mát rượi để trở về mái ấm ký túc xá ở phố Anwa.

Tô và Ki tung tăng tiễn chân đến tận cổng sắt, thỉnh thoảng lại sủa vang vài tiếng mừng rỡ, đuôi ngoáy tít mù nhìn theo bóng xe của cả nhà khuất dần sau khúc quanh đại lộ ngoại ô. Dưới ánh hoàng hôn nhập nhoạng, hai bóng cún trắng in rõ bên cánh cổng sắt, bình yên như một bức tranh tuyệt mỹ.

Cả Suzuno, Naomi và đại gia đình Hikawa đều không thể ngờ rằng\dots

Đây chính là buổi chiều yên ả cuối cùng mà Tô và Ki còn được tự do chạy nhảy dưới hình hài của hai chú chó bốn chân.

\ct

Đêm muộn.

Gió đêm mùa hạ thổi ràn rạt qua bãi đất hoang vắng ngắt. Cả khu vực ngoại ô Kawachi chìm vào giấc ngủ sâu thẳm, chỉ còn tiếng côn trùng rả rích trong những bụi cỏ dại rậm rạp và tiếng cành cây rẻ quạt khẽ lay động trong bóng tối tịch mịch.

Xưởng cơ khí đã tắt hết đèn, cánh cổng sắt nặng nề khóa chặt. Phòng thí nghiệm của Suzuno đứng lặng lẽ bên bức tường rào tôn gỉ sét, ô cửa thông gió phản chiếu ánh trăng mờ nhạt lạnh lẽo.

``Xột xoạt\dots''

Một âm thanh kỳ lạ vang lên từ phía hàng rào ngăn cách bãi đất hoang với con hẻm tối ngoài phố lớn.

Bụi cây rậm rạp khẽ rẽ sang hai bên. Một bóng người nhỏ nhắn, khoác chiếc áo blouse dài màu hồng phấn trườn nhẹ qua khe hở của tấm tôn rào, tiếp đất bằng những bước chân nhẹ tênh không một tiếng động.

Dưới ánh trăng bàng bạc chiếu rọi qua kẽ lá, một đôi tai thú màu hồng tro dựng đứng trên đỉnh đầu khẽ giật giật nhạy bén, liên tục quét các hướng để lắng nghe từng dao động âm thanh trong bán kính hàng trăm mét. Đôi mắt sau cặp kính bảo hộ gọng tròn lóe lên những tia sáng xanh ma mị, sắc sảo và đầy vẻ tính toán của một nhà khoa học thiên tài điên rồ.

Chiếc đuôi sói phía sau khẽ đung đưa theo từng nhịp bước chân lén lút.

Bóng người ấy dừng lại trước mép hiên phòng thí nghiệm của Suzuno. Bàn tay đeo găng cao su màu đen khẽ chạm vào chiếc khay nhôm đặt trên bàn thao tác ngoài trờ, nơi cuộn lông tơ trắng dính tóc của nhà Hikawa vẫn đang nằm nguyên ở đó.

Một nụ cười tinh quái, sắc lạnh như lưỡi dao mổ khẽ nở trên khóe môi: ``Kufufu\dots\ Mẫu vật hữu cơ mang cấu trúc đa tổ hợp gen cực kỳ thú vị\dots\ Đúng là không uổng công ta lặn lội đến tận xưởng cơ khí hẻo lánh này.''

Cô nàng cẩn thận dùng chiếc kẹp chuyên dụng nhặt cuộn lông tơ bỏ vào ống nghiệm vô trùng, đậy nắp kín rồi cất vào túi áo blouse.

Sau đó, đôi mắt màu ngọc bích sau cặp kính bảo hộ chầm chậm xoay hướng, nhìn thẳng qua khoảng sân tối tăm, khóa chặt mục tiêu vào tảng đá hoa cương trắng nằm dưới gốc cây rẻ quạt cổ thụ, nơi an nghỉ của cô chó Mi.

Đôi tai sói khẽ vểnh cao, phát ra một tiếng cười thì thầm đầy đắc thắng trong màn đêm tĩnh mịch:

``Hóa ra `vật chủ' nguyên bản của tổ hợp gen này được chôn ở đằng kia sao? Được rồi\dots\ Đến lúc đánh thức ngươi dậy rồi, người bạn nhỏ.''

Bóng đen lướt nhanh qua bãi cỏ trong đêm, hướng thẳng về phía ngôi mộ mới dưới gốc cây rẻ quạt, mở màn cho một biến cố chấn động sắp sửa giáng xuống gia đình Hikawa\dots

\end{document}